\documentclass[a4paper,11pt]{article}
\usepackage[T1]{fontenc}
\usepackage[utf8]{inputenc}
\usepackage[ngerman]{babel}
\usepackage{lmodern}
\usepackage{amsmath,amssymb}
\usepackage{geometry}
\usepackage[table]{xcolor}
\usepackage{tikz}
\usepackage{eso-pic}
\geometry{paper=a4paper,top=0cm,bottom=0cm,left=0cm,right=0cm,headheight=0pt,headsep=0pt,footskip=0pt}
\pagestyle{empty}
\definecolor{examblue}{RGB}{0,70,130}
\definecolor{examgreen}{RGB}{0,110,60}
\definecolor{cardgray}{RGB}{248,248,248}
\newlength{\cardw}\newlength{\cardh}
\setlength{\cardw}{9.80cm}\setlength{\cardh}{5.24cm}
\AddToShipoutPictureFG{\begin{tikzpicture}[remember picture,overlay]
\draw[gray!70,densely dotted,line width=0.7pt]([xshift=10.5cm]current page.south west)--([xshift=10.5cm]current page.north west);
\foreach \y in {5.94,11.88,17.82,23.76}{\draw[gray!70,densely dotted,line width=0.7pt]([yshift=\y cm]current page.south west)--([yshift=\y cm]current page.south east);}
\end{tikzpicture}}
\newcommand{\checkfeld}{\raisebox{-0.15ex}{\fbox{\rule{0pt}{1.15ex}\rule{1.15ex}{0pt}}}}
\newcommand{\frontcard}[3]{\begin{tikzpicture}
\node[draw=examblue,line width=0.8pt,rounded corners=4pt,fill=white,minimum width=\cardw,minimum height=\cardh,text width=8.75cm,align=center,inner sep=8pt](card){\Large\bfseries #3};
\node[anchor=north west,fill=examblue,text=white,rounded corners=3pt,font=\bfseries\small,inner xsep=7pt,inner ysep=4pt]at([xshift=7pt,yshift=-7pt]card.north west){Karte #1};
\node[anchor=south,text=examblue,font=\scriptsize\bfseries]at([yshift=21pt]card.south){NaWi $\cdot$ #2};
\node[anchor=south,text=examblue,font=\scriptsize]at([yshift=6pt]card.south){Gewusst:\enspace\checkfeld\enspace\checkfeld\enspace\checkfeld\enspace\checkfeld\enspace\checkfeld};
\end{tikzpicture}}
\newcommand{\backcard}[3]{\begin{tikzpicture}
\node[draw=examgreen,line width=0.8pt,rounded corners=4pt,fill=cardgray,minimum width=\cardw,minimum height=\cardh,text width=8.75cm,align=center,inner sep=9pt](card){\large #3};
\node[anchor=north west,fill=examgreen,text=white,rounded corners=3pt,font=\bfseries\small,inner xsep=7pt,inner ysep=4pt]at([xshift=7pt,yshift=-7pt]card.north west){Karte #1};
\node[anchor=south,text=examgreen,font=\scriptsize\bfseries]at([yshift=21pt]card.south){NaWi $\cdot$ #2};
\node[anchor=south,text=examgreen,font=\scriptsize]at([yshift=6pt]card.south){Gewusst:\enspace\checkfeld\enspace\checkfeld\enspace\checkfeld\enspace\checkfeld\enspace\checkfeld};
\end{tikzpicture}}
\newcommand{\blankcard}{\begin{tikzpicture}\node[minimum width=\cardw,minimum height=\cardh]{};\end{tikzpicture}}
\newcommand{\cardcell}[1]{\parbox[c][5.94cm][c]{10.50cm}{\centering #1}}
\newcommand{\cardrow}[2]{\noindent\cardcell{#1}\cardcell{#2}\par}
\begin{document}
\setlength{\topskip}{0pt}
\offinterlineskip
% Karten 188 bis 217: Vorder- und Rueckseiten
\cardrow
{\frontcard{188}{Kernphysik}{Beschreibe den Aufbau eines Atoms.}}
{\frontcard{189}{Kernphysik}{Beschreibe den Aufbau des Atomkerns.}}
\cardrow
{\frontcard{190}{Kernphysik}{Was versteht man unter einem Isotop?}}
{\frontcard{191}{Kernphysik}{Welcher Kernbaustein ist für ein Element immer gleich?}}
\cardrow
{\frontcard{192}{Kernphysik}{Beschreibe die Kernkraft in ihren wesentlichen Punkten.}}
{\frontcard{193}{Kernphysik}{Was ist der Massendefekt?}}
\cardrow
{\frontcard{194}{Kernphysik}{Wie wird Energie bei der Kernspaltung gewonnen?}}
{\frontcard{195}{Kernphysik}{Wie wird Energie bei der Kernfusion gewonnen?}}
\cardrow
{\frontcard{196}{Kernphysik}{Was ist Alphastrahlung?}}
{\frontcard{197}{Kernphysik}{Was ist Gammastrahlung?}}
\newpage
\cardrow
{\backcard{189}{Kernphysik}{Der Kern besteht aus Protonen und Neutronen, den Nukleonen, die durch die starke Kernkraft gebunden sind.}}
{\backcard{188}{Kernphysik}{Ein Atom besteht aus kleinem Kern mit Protonen und Neutronen sowie einer Elektronenhülle; neutral gilt Protonenzahl gleich Elektronenzahl.}}
\cardrow
{\backcard{191}{Kernphysik}{Die Protonenzahl beziehungsweise Ordnungszahl.}}
{\backcard{190}{Kernphysik}{Atome desselben Elements mit gleicher Protonenzahl, aber unterschiedlicher Neutronen- und Massenzahl.}}
\cardrow
{\backcard{193}{Kernphysik}{Die Kernmasse ist kleiner als die Summe der freien Nukleonenmassen; \(E=\Delta m c^2\) entspricht der Bindungsenergie.}}
{\backcard{192}{Kernphysik}{Sie wirkt zwischen Nukleonen, ist sehr stark und kurzreichweitig, meist anziehend und bei extrem kleinen Abständen abstoßend.}}
\cardrow
{\backcard{195}{Kernphysik}{Leichte Kerne verschmelzen; der Massendefekt wird gemäß \(E=\Delta m c^2\) als Energie frei.}}
{\backcard{194}{Kernphysik}{Ein schwerer Kern wird gespalten; der Massendefekt wird gemäß \(E=\Delta m c^2\) als Energie frei.}}
\cardrow
{\backcard{197}{Kernphysik}{Hochenergetische elektromagnetische Strahlung aus angeregten Kernen; ungeladen und sehr durchdringend.}}
{\backcard{196}{Kernphysik}{Heliumkerne aus zwei Protonen und zwei Neutronen; zweifach positiv, stark ionisierend und wenig durchdringend.}}
\newpage
\cardrow
{\frontcard{198}{Kernphysik}{Was ist Betastrahlung?}}
{\frontcard{199}{Kernphysik}{Gib die allgemeine Zerfallsgleichung für Gammastrahlung an.}}
\cardrow
{\frontcard{200}{Kernphysik}{Gib die allgemeine Zerfallsgleichung für Alphastrahlung an.}}
{\frontcard{201}{Kernphysik}{Gib die allgemeine Zerfallsgleichung für den Beta-plus-Zerfall an.}}
\cardrow
{\frontcard{202}{Kernphysik}{Was versteht man unter einer Zerfallsreihe?}}
{\frontcard{203}{Kernphysik}{Wie lautet das Zerfallsgesetz?}}
\cardrow
{\frontcard{204}{Kernphysik}{Beschreibe die Größen- und Massenverhältnisse in einem Atom.}}
{\frontcard{205}{Kernphysik}{Gib die allgemeine Zerfallsgleichung für den Beta-minus-Zerfall an.}}
\cardrow
{\frontcard{206}{Kernphysik}{Definiere die physikalische Größe Aktivität. Gib auch Formelzeichen und SI-Einheit an.}}
{\frontcard{207}{Kernphysik}{Definiere die physikalische Größe Energiedosis. Gib auch Formelzeichen und SI-Einheit an.}}
\newpage
\cardrow
{\backcard{199}{Kernphysik}{\({}^{A}_{Z}X^*\rightarrow{}^{A}_{Z}X+\gamma\).}}
{\backcard{198}{Kernphysik}{Beim \(\beta^-\)-Zerfall entstehen Elektron und Antineutrino, beim \(\beta^+\)-Zerfall Positron und Neutrino.}}
\cardrow
{\backcard{201}{Kernphysik}{\({}^{A}_{Z}X\rightarrow{}^{A}_{Z-1}Y+e^++\nu_e\).}}
{\backcard{200}{Kernphysik}{\({}^{A}_{Z}X\rightarrow{}^{A-4}_{Z-2}Y+{}^{4}_{2}\mathrm{He}\).}}
\cardrow
{\backcard{203}{Kernphysik}{\(N(t)=N_0e^{-\lambda t}\), außerdem \(T_{1/2}=\ln2/\lambda\).}}
{\backcard{202}{Kernphysik}{Eine Folge radioaktiver Zerfälle über instabile Tochterkerne bis zu einem stabilen Endkern.}}
\cardrow
{\backcard{205}{Kernphysik}{\({}^{A}_{Z}X\rightarrow{}^{A}_{Z+1}Y+e^-+\overline{\nu}_e\).}}
{\backcard{204}{Kernphysik}{Atom etwa \(10^{-10}\,\mathrm m\), Kern etwa \(10^{-15}\,\mathrm m\); fast die ganze Masse liegt im Kern.}}
\cardrow
{\backcard{207}{Kernphysik}{Absorbierte Energie pro Masse: \(D=E_\mathrm{abs}/m\); SI-Einheit Gray, \(1\,\mathrm{Gy}=1\,\mathrm{J/kg}\).}}
{\backcard{206}{Kernphysik}{Zerfälle pro Zeit: \(A=-\mathrm dN/\mathrm dt=\lambda N\); SI-Einheit Becquerel, \(1\,\mathrm{Bq}=1\,\mathrm{s^{-1}}\).}}
\newpage
\cardrow
{\frontcard{208}{Kernphysik}{Definiere die physikalische Größe Äquivalentdosis. Gib auch Formelzeichen und SI-Einheit an.}}
{\frontcard{209}{Kernphysik}{Was versteht man unter ionisierender Strahlung?}}
\cardrow
{\frontcard{210}{Kernphysik}{Beschreibe in den wesentlichen Punkten die biologische Strahlenwirkung.}}
{\frontcard{211}{Kernphysik}{Wie kann man sich vor radioaktiver Strahlung schützen?}}
\cardrow
{\frontcard{212}{Kernphysik}{Benenne drei Hauptkomponenten eines typischen thermischen Reaktorkerns.}}
{\frontcard{213}{Kernphysik}{Was ist die Aufgabe des Moderators?}}
\cardrow
{\frontcard{214}{Kernphysik}{Was ist die Aufgabe eines Regelstabes?}}
{\frontcard{215}{Kernphysik}{Was versteht man unter der kritischen Masse?}}
\cardrow
{\frontcard{216}{Kernphysik}{Was passiert bei der Kernfusion?}}
{\frontcard{217}{Kernphysik}{Unter welchen Bedingungen ist Kernfusion möglich?}}
\newpage
\cardrow
{\backcard{209}{Kernphysik}{Strahlung mit genügend Energie, Elektronen aus Atomen oder Molekülen zu entfernen und Ionen zu erzeugen.}}
{\backcard{208}{Kernphysik}{Biologisch gewichtete Dosis: \(H=\sum_R w_R D_R\); SI-Einheit Sievert.}}
\cardrow
{\backcard{211}{Kernphysik}{ALARA: Zeit kurz, Abstand groß, geeignete Abschirmung; außerdem Kontamination vermeiden und Dosimeter verwenden.}}
{\backcard{210}{Kernphysik}{Sie kann Moleküle und DNA direkt oder über Radikale schädigen; Folgen hängen von Dosis, Strahlungsart, Gewebe und Zeitverlauf ab.}}
\cardrow
{\backcard{213}{Kernphysik}{Er bremst schnelle Neutronen ab, ohne sie möglichst stark zu absorbieren.}}
{\backcard{212}{Kernphysik}{Brennelemente, Moderator und Regelstäbe; zusätzlich führt ein Kühlsystem Wärme ab.}}
\cardrow
{\backcard{215}{Kernphysik}{Kleinste Materialmenge, bei der unter gegebenen Bedingungen eine selbsterhaltende Kettenreaktion möglich ist.}}
{\backcard{214}{Kernphysik}{Er absorbiert Neutronen und regelt dadurch Leistung beziehungsweise Abschaltung der Kettenreaktion.}}
\cardrow
{\backcard{217}{Kernphysik}{Extrem heißes Plasma, genügend hohe Teilchendichte und ausreichend lange Einschlusszeit.}}
{\backcard{216}{Kernphysik}{Zwei leichte Kerne verschmelzen zu einem schwereren; ein Teil der Masse wird in Energie umgewandelt.}}
\end{document}
